\FloatBarrier
\section{本章小结}
\label{sec:subset-selection-chapter-summary}

本章回答了名额或经费有限时怎样选子集。两种贪心有相同的骨架：从空集出发，逐轮维护已选集合与当前值，相对于当前集合重算增益，加入一个元素。基数约束下比较的是边际增益本身，预算约束下比较的是单位成本的增益，并且只在付得起的元素中比较。两者都只需多项式次求值，分别至多$nk+1$次和$n(n+1)/2+1$次，总代价还要乘以每次求值的开销。

贪心规则本身对任何集合函数都能执行，但结果有没有保证取决于目标性质和约束。基数约束下最大化归一化、单调、次模的函数时，最优集合的至多$k$个元素在每一轮都仍然可选；由次模性，当前最好的元素至少补上剩余差距的$1/k$，差距按几何级数缩小，贪心值不低于最优值的$1-(1-1/k)^k$倍，从而不低于$1-e^{-1}$倍。撤去次模性，贪心看不到只有同时选中才有价值的互补组合；撤去单调性，强制选满会降低目标。两种情形下比例保证都不再成立。

惰性贪心是同一规则的另一种实现。次模性使旧增益成为当前增益的上界，优先队列据此只重算排在最前的元素；在一致的平局规则下，它与普通贪心逐轮选出相同的元素，所以继承同样的保证。它的正确性只需要次模性，节省多少求值没有一般的保证，最坏情形与普通贪心相同。

预算约束改变了可行性，也改变了分析。单位成本规则在传感器实例中找到最优解，却在成本与价值分别为$(1,2)$与$(100,100)$的例子中可以任意差：选进便宜元素以后，最优集合的元素可能再也放不进剩余预算，差距论证中“最优集合的元素仍可选”这一步失效。最佳单元素比较能修复这个例子，但本章没有证明它的一般比例。使用任何一种选择方法之前，都要写明目标方向、约束形式、函数是否归一化、单调和次模。

下一章转向最小化方向，没有基数或预算约束，但要求目标由一元项和成对项组成；在相应的次模条件下，最小割给出精确解。
